% ============================================================
%  Introduction Générale — Mémoire de Master BIBDA
%  CHOUFA Zouhair
% ============================================================
\chapter*{Introduction Générale}
\addcontentsline{toc}{chapter}{Introduction Générale}
\label{chap:introduction_generale}
 
\section*{Contexte général}
\addcontentsline{toc}{section}{Contexte général}
 
L'ère du Big Data a profondément transformé la manière dont les
organisations produisent, stockent et exploitent leurs données.
Parmi les cinq dimensions classiquement associées à ce paradigme —
Volume, Vélocité, Variété, Véracité et Valeur —, la \textbf{Variété}
constitue sans doute le défi le plus structurel pour les entreprises
modernes. Une organisation typique dispose rarement d'une source de
données unique et homogène~: ses informations sont dispersées entre
systèmes CRM, bases de données transactionnelles, fichiers plats
CSV, systèmes d'information des ressources humaines (SIRH) ou
tableurs, chacun obéissant à des conventions de nommage, des schémas
et des formats qui lui sont propres. Cette fragmentation, souvent
qualifiée de phénomène des \textit{silos de données}, constitue un
frein majeur à toute exploitation analytique transversale et
cohérente de l'information disponible.
 
Les \textbf{Knowledge Graphs} (KG), fondés sur le modèle de données
RDF (\textit{Resource Description Framework}) et sur les standards du
Web sémantique — OWL, SPARQL, SHACL —, se sont imposés au cours de la
dernière décennie comme un cadre formel particulièrement adapté pour
répondre à ce défi d'intégration. En représentant les entités et
leurs relations sous forme de triplets sujet-prédicat-objet ancrés
dans une ontologie de domaine partagée, les KG permettent de dépasser
la rigidité du modèle relationnel classique et d'unifier
sémantiquement des sources hétérogènes autour d'un vocabulaire
commun. Cette approche a démontré sa valeur dans des contextes aussi
variés que les moteurs de recherche, les systèmes de recommandation,
la biomédecine ou l'intégration de données d'entreprise.
 
\section*{Problématique}
\addcontentsline{toc}{section}{Problématique}
 
Si la valeur des Knowledge Graphs pour l'intégration de données
n'est plus à démontrer, leur \textbf{construction automatisée} à
partir de sources tabulaires hétérogènes demeure un problème ouvert
et largement non résolu de façon industrialisable. La littérature
existante propose des solutions ponctuelles pour chacune des étapes
du processus — schema matching, triplification déclarative via les
standards R2RML/RML, résolution d'entités, validation de la
qualité du graphe produit — mais rarement de manière intégrée au
sein d'un pipeline unique, reproductible et évaluable de bout en
bout. De plus, l'émergence récente des grands modèles de langage
(\textit{Large Language Models}, LLM) ouvre des perspectives
prometteuses pour automatiser des tâches historiquement coûteuses en
intervention humaine, telles que l'alignement sémantique de schémas
hétérogènes, mais leur intégration dans des architectures de
production robustes, maîtrisées en coût et rigoureusement évaluées,
reste largement à explorer.
 
Ce mémoire s'inscrit directement dans ce contexte et cherche à
répondre à la question de recherche suivante~: \textit{comment
concevoir un pipeline automatisé, modulaire et reproductible, capable
de transformer des données tabulaires hétérogènes et multi-sources en
un Knowledge Graph unifié, tout en garantissant la qualité et la
validité formelle du graphe produit à chaque étape du processus~?}
 
\section*{Objectifs du mémoire}
\addcontentsline{toc}{section}{Objectifs du mémoire}
 
Pour répondre à cette problématique, ce travail poursuit quatre
objectifs principaux~:
 
\begin{itemize}
    \item \textbf{Concevoir} une architecture de pipeline modulaire
    couvrant l'ensemble de la chaîne de transformation~: ingestion et
    profilage qualité des sources, alignement sémantique des schémas,
    triplification RDF déclarative et résolution d'entités
    inter-sources ;
    \item \textbf{Combiner} de manière hybride des approches
    symboliques (embeddings Sentence-BERT) et neuronales fondées sur
    les LLM (agent LangChain) pour le \textit{schema matching}, en
    visant à maximiser la qualité d'alignement tout en maîtrisant le
    coût computationnel des appels aux modèles de langage ;
    \item \textbf{Automatiser} la génération des règles de mapping
    YARRRML et leur exécution via le moteur Morph-KGC, ainsi que la
    validation formelle du graphe de connaissances produit au moyen
    de contraintes SHACL ;
    \item \textbf{Évaluer} quantitativement chaque module du pipeline
    sur un gold standard annoté, et discuter de manière critique la
    validité, les limites et la généralisabilité des résultats
    obtenus.
\end{itemize}
 
\section*{Contributions}
\addcontentsline{toc}{section}{Contributions}
 
La contribution de ce mémoire se situe à l'intersection de trois
axes de recherche généralement traités séparément dans la
littérature~: la triplification déclarative de données (RML/YARRRML),
le \textit{schema matching} hybride, et l'automatisation par LLM.
Le pipeline proposé intègre ces trois dimensions au sein d'une
architecture unique, entièrement conteneurisée via Docker Compose,
et dont chaque module fait l'objet d'une évaluation quantitative
indépendante sur un gold standard construit spécifiquement pour ce
travail. Cette approche intégrée et évaluable de bout en bout
constitue, à notre connaissance, une réponse originale aux lacunes
identifiées dans l'état de l'art~: fragmentation des pipelines
existants, absence d'évaluation standardisée, reproductibilité
insuffisante et déficit d'études portant sur la scalabilité.
 
\section*{Organisation du mémoire}
\addcontentsline{toc}{section}{Organisation du mémoire}
 
Ce mémoire est structuré en quatre chapitres. Le
\textbf{Chapitre~1} dresse un état de l'art structuré autour du
paradigme Big Data, des fondements des Knowledge Graphs, des
approches de transformation tabulaire vers RDF et de
\textit{schema matching}, ainsi que de l'apport récent des LLM à ces
problématiques ; il se conclut par le positionnement précis de notre
contribution par rapport aux travaux existants. Le
\textbf{Chapitre~2} présente la méthodologie et la conception
détaillée de l'architecture du pipeline proposé, module par module,
ainsi que le protocole d'évaluation et la constitution du gold
standard. Le \textbf{Chapitre~3} décrit l'implémentation concrète de
ce pipeline dans un environnement Cloud réel, les résultats obtenus
pour chaque module, ainsi qu'une étude d'ablation isolant la
contribution propre de chaque composant du module de
\textit{schema matching}. Enfin, le \textbf{Chapitre~4} propose une
évaluation approfondie et comparative des résultats~: calibration
des hyperparamètres critiques du pipeline, positionnement par rapport
à l'état de l'art, et discussion critique structurée autour des
menaces classiques à la validité expérimentale. Une conclusion
générale synthétise les apports de ce travail, en discute les limites
et propose des perspectives de recherche future.
 